% TOP_PROBLEM: {"id": 4500001, "title": "Periods of Pseudo-Integrable Billiards", "category_id": 11, "category": {"id": 11, "name": "dynamical_systems", "display_name": "Dynamical Systems", "description": "Problems about long-term behavior of deterministic systems, Hamiltonian dynamics, and periodic orbits.", "slug": "dynamical-systems", "order_index": 11, "created_at": "2026-07-31T15:26:25.671Z"}, "status": "open"}
% FIRST_POSTED: null
% ENTRY_KIND: statement_only
% Imported from: batch14.tex; UnsolvedMath ID 4500001
% Compile twice: lualatex 4500001.tex
\documentclass[11pt]{article}
\usepackage{fontspec}
\setmainfont{Latin Modern Roman}
\usepackage{amsmath,amssymb,mathtools,unicode-math}
\setmathfont{Latin Modern Math}
\usepackage{xurl,hyperref}
\hypersetup{hidelinks,pdftitle={UnsolvedMath Problem 4500001}}
\newcommand{\sref}[2]{\hyperlink{source:#1:#2}{[S#2]}}
\newcommand{\cataloguescope}[1]{\paragraph{Mathematical question.}#1}
\begin{document}
% UnsolvedMath ID 4500001; source catalogue code AMR-044-0001
\setcounter{section}{4500000}
\section{Periods of Pseudo-Integrable Billiards}\label{4500001}
{\small\sffamily UnsolvedMath ID 4500001 · Dynamical Systems}
\subsection{Definitions and mathematical statement}

\paragraph{Shared notation.}$\mathbb N=\{1,2,\ldots\}$, and $\mathbb Z,\mathbb Q,\mathbb R,\mathbb C$ denote the integers, rationals, reals and complex numbers. Any other conventions are specified locally.


Choose radii \(R_1>R_2>r>0\). In the Euclidean plane take the billiard table whose boundary consists of
\[
\Gamma_1=\{(x,y):x^2+y^2=R_1^2,\ x\ge0\},\qquad
\Gamma_2=\{(x,y):x^2+y^2=R_2^2,\ x\le0\},
\]
and the segments \(\{(0,y):R_2\le y\le R_1\}\) and \(\{(0,y):-R_1\le y\le-R_2\}\). The billiard particle moves along straight segments at constant speed and reflects elastically, with equal incidence and reflection angles, at smooth boundary points. At a convex right-angle vertex the limiting reflection reverses the velocity and counts as two bounces. A trajectory that reaches a reflex-angle vertex terminates; complete trajectories avoid those vertices.
Restrict attention to complete trajectories whose segments are tangent to the concentric circle \(x^2+y^2=r^2\), called their caustic. For a circular table of radius \(R_i\) with this caustic, the central angle between consecutive impacts is \(2\arccos(r/R_i)\); the rotation number is
\[\rho_i=\frac1\pi\arccos\left(\frac r{R_i}\right),\qquad
0<\rho_2<\rho_1<\frac12.\]
A trajectory is periodic when it returns to the same position and velocity after finitely many reflections; its minimal bounce period is the least positive number of bounces required for such a return, with the vertex convention above.
\paragraph{Statement recorded in the supplied source.}
For every admissible pair \((\rho_1,\rho_2)\), determine which complete trajectories with the fixed circular caustic are periodic and which minimal bounce periods occur. In particular, describe the dependence of periodicity and possible periods on the arithmetic of the two rotation numbers. \sref{4500001}{2}
\subsection{Short English statement}
For the table made from two concentric semicircles and two joining segments, classify periodic trajectories tangent to a fixed concentric circle and determine all their possible periods from the two circular rotation numbers.
\subsection{Sources}
\begin{description}
\item[\hypertarget{source:4500001:1}{S1}] ULAM.AI and the UnsolvedMath Contributors, UnsolvedMath: A Curated Collection of Open Mathematics Problems, record 4500001 (AMR-044-0001). CC BY 4.0. \url{https://huggingface.co/datasets/ulamai/UnsolvedMath}.
\item[\hypertarget{source:4500001:2}{S2}] Vladimir Dragović and Milena Radnović, Pseudo-integrable Billiards and Arithmetic Dynamics (2012), arXiv:1206.0163, Section 1, vertex conventions; Section 3, domain D0; Section 4, circular examples; Section 7, Theorem 7.1 and rotation-dependent interval exchanges. \url{https://arxiv.org/pdf/1206.0163}.
\end{description}
\label{end:4500001}
\end{document}
